\documentclass[12pt]{article}
\usepackage{graphicx} % Required for inserting images

\title{S2 G3 ITS L5 Scribe CoT AU2440248 Aryan}
\author{Aryan Thakkar}
\date{\today}

\begin{document}

\maketitle
\section*{CSE 400: Fundamentals of Probability in Computing}
\subsection*{Lecture 5: Bayes' Theorem, Random Variables, and Probability Mass Function}

\section{Bayes' Theorem}

\subsection{Weighted Average of Conditional Probabilities}

Let $A$ and $B$ be events.

We may express:
\[
A = AB \cup AB^c
\]

Reason: For an outcome to be in $A$, it must either:
\begin{itemize}
    \item be in both $A$ and $B$, or
    \item be in $A$ but not in $B$
\end{itemize}

Since $AB$ and $AB^c$ are mutually exclusive, by Axiom 3:
\[
\Pr(A) = \Pr(AB) + \Pr(AB^c)
\]

Using conditional probability:
\[
\Pr(AB) = \Pr(A \mid B)\Pr(B)
\]
\[
\Pr(AB^c) = \Pr(A \mid B^c)\Pr(B^c)
\]

Thus:
\[
\Pr(A) = \Pr(A \mid B)\Pr(B) + \Pr(A \mid B^c)[1-\Pr(B)]
\]

\textbf{Key Statement:}

The probability of event $A$ is a weighted average of the conditional probabilities, with weights given as the probability of the event on which it is conditioned.

\subsection{Learning by Example}

\subsubsection*{Example 3.1 (Part 1/2)}

An insurance company divides people into two classes:
\begin{itemize}
    \item Accident-prone
    \item Not accident-prone
\end{itemize}

Statistics:
\begin{itemize}
    \item Accident-prone person has an accident within 1 year with probability $0.4$
    \item Not accident-prone person has an accident within 1 year with probability $0.2$
\end{itemize}

Assume:
\[
\Pr(A)=0.3
\]
where $A$ is the event that the policyholder is accident prone.

Let:
\[
A_1 = \mathrm{event that the policyholder has an accident within 1 year}
\]

We want:
\[
\Pr(A_1)
\]

Solution:

Condition on whether the policyholder is accident prone:
\[
\Pr(A_1)=\Pr(A_1\mid A)\Pr(A)+\Pr(A_1\mid A^c)\Pr(A^c)
\]

Substitute values:
\[
\Pr(A_1)=(0.4)(0.3)+(0.2)(0.7)
\]

So:
\[
\Pr(A_1)=0.12+0.14=0.26
\]

Thus:
\[
\Pr(A_1)=0.26
\]

\subsubsection*{Example 3.1 (Part 2/2)}

Suppose a new policyholder has an accident within a year.

What is:
\[
\Pr(A\mid A_1)?
\]

Solution:

\[
\Pr(A\mid A_1)=\frac{\Pr(AA_1)}{\Pr(A_1)}
\]

\[
=\frac{\Pr(A)\Pr(A_1\mid A)}{0.26}
\]

\[
=\frac{(0.3)(0.4)}{0.26}
=\frac{0.12}{0.26}
=\frac{6}{13}
\]

Thus:
\[
\Pr(A\mid A_1)=\frac{6}{13}
\]

\subsection{Law of Total Probability and Bayes Formula}

\subsubsection*{Law of Total Probability (Formula 3.4)}

Suppose $B_1,B_2,\dots,B_n$ are mutually exclusive events such that:
\[
B=\bigcup_{i=1}^{n}B_i
\]

Then:
\[
\Pr(A)=\sum_{i=1}^{n}\Pr(AB_i)
\]

We obtain:
\[
\Pr(A)=\sum_{i=1}^{n}\Pr(A\mid B_i)\Pr(B_i)
\]

This is the Law of Total Probability.

\subsubsection*{Bayes Formula (Proposition 3.1)}

Using:
\[
\Pr(AB_i)=\Pr(B_i\mid A)\Pr(A)
\]
We get:
\[
\Pr(B_i\mid A)=
\frac{\Pr(A\mid B_i)\Pr(B_i)}
{\sum_{j=1}^{n}\Pr(A\mid B_j)\Pr(B_j)}
\]
This is the Bayes Formula.
\textbf{Definitions:}
\begin{itemize}
    \item $\Pr(B_i)$ is the \textit{apriori probability}
    \item $\Pr(B_i\mid A)$ is the \textit{posteriori probability}
\end{itemize}
\subsection{Learning by Example}
\subsubsection*{Example 3.2 (Cards Problem)}
There are 3 cards:
\begin{enumerate}
    \item Red--Red
    \item Black--Black
    \item Red--Black
\end{enumerate}
One card is randomly selected and placed down.
Given: Upper side is red.
Find probability that the other side is black.
Let:
\[
RR=\text{chosen card is all red}
\]
\[
BB=\text{chosen card is all black}
\]
\[
RB=\text{chosen card is red--black}
\]
\[
R=\text{event that upturned side is red}
\]
We want:
\[
\Pr(RB\mid R)
\]
Substitute:
\[
\Pr(RB\mid R)=\frac{\Pr(RB\cap R)}{\Pr(R)}
\]
\[
=\frac{\Pr(R\mid RB)\Pr(RB)}
{\Pr(R\mid RR)\Pr(RR)+\Pr(R\mid RB)\Pr(RB)+\Pr(R\mid BB)\Pr(BB)}
\]
Now:
\[
\Pr(R\mid RR)=1,\quad \Pr(R\mid RB)=\frac12,\quad \Pr(R\mid BB)=0
\]
Thus:
\[
\Pr(RB\mid R)=
\frac{(1/2)(1/3)}
{(1)(1/3)+(1/2)(1/3)+(0)(1/3)}
\]
\[
=\frac{1/6}{1/3+1/6}
=\frac{1/6}{1/2}
=\frac{1}{3}
\]
Thus:
\[
\Pr(RB\mid R)=\frac13
\]
\section{Random Variables}
\subsection{Motivation and Concept}
When an experiment is performed, we are often interested in some function of the outcome.
Examples:
\begin{itemize}
    \item Dice tossing: focus on the sum
    \item Coin flipping: focus on number of heads
\end{itemize}
These real-valued functions are called Random Variables.
\subsubsection*{Definition}
A random variable on sample space $\Omega$ is a function:
\[
X:\Omega\to\mathbb{R}
\]
assigning each sample point $\omega\in\Omega$ a real number $X(\omega)$.
\subsubsection*{Distribution}
Important aspects:
\begin{itemize}
    \item The set of values it can take
    \item The probabilities with which it takes those values
\end{itemize}
Event notation:
\[
\{\omega\in\Omega:X(\omega)=a\}
\]
Abbreviated as:
\[
X=a
\]
Probability:
\[
\Pr[X=a]
\]
The collection of these probabilities is called the distribution of $X$.
\subsection{Random Variable Example}
\subsubsection*{Example 1: Tossing 3 Fair Coins}
Let $Y$ denote number of heads.
Possible values:
\[
0,1,2,3
\]
Probabilities:
\[
P(Y=0)=P(t,t,t)=\frac18
\]
\[
P(Y=1)=\frac38
\]
\[
P(Y=2)=\frac38
\]
\[
P(Y=3)=\frac18
\]
Since $Y$ must take one of these values:
\[
1=\sum_{i=0}^{3}P(Y=i)
\]
\section{Probability Mass Function (PMF)}
\subsection{Concept}
A random variable that can take on at most a countable number of values is said to be discrete.

Let $X$ be discrete with range:
\[
R_X=x_1,x_2,x_3,\dots
\]
The function:
\[
p(x_k)=\Pr(X=x_k)
\]
is called the Probability Mass Function (PMF) of $X$.
Since $X$ must take one of the values:
\[
\sum_k p(x_k)=1
\]
\subsection{PMF Example}
The PMF is given by:
\[
p(i)=c\frac{\lambda^i}{i!},\quad i=0,1,2,\dots
\]
where $\lambda>0$.
Find:
\[
P(X=0)
\quad\text{and}\quad
P(X>2)
\]
Solution:
Since:
\[
\sum_{i=0}^{\infty}p(i)=1
\]
But:
\[
\sum_{i=0}^{\infty}c\frac{\lambda^i}{i!}=1
\]
So:
\[
c\sum_{i=0}^{\infty}\frac{\lambda^i}{i!}=1
\]
\[
ce^\lambda=1
\]

Thus:
\[
c=e^{-\lambda}
\]
\subsubsection*{(1) Finding $P(X=0)$}
\[
P(X=0)=p(0)=c\frac{\lambda^0}{0!}=e^{-\lambda}
\]
\subsubsection*{(2) Finding $P(X>2)$}
\[
P(X>2)=1-P(X\le 2)
\]
\[
=1-[P(X=0)+P(X=1)+P(X=2)]
\]
\section*{Lecture 5 Complete}
\end{document}
